%%%%%%%%%%%%%%%
% This CV example/template is based on my own
% CV which I (lamely attempted) to clean up, so that
% it's less of an eyesore and easier for others to use.
%
% LianTze Lim (liantze@gmail.com)
% 23 Oct, 2022
% 24 Aug, 2024 -- Updated X (Twitter) icon

% Zhiqiang Ma (zhiqiang.ma97@gmail.com)
% 2025/5/20, Revised
%% 编译方式 (Compilation):
%% 1. xelatex cv-main
%% 2. biber cv-main
%% 3. xelatex cv-main
%% 4. xelatex cv-main

\documentclass[a4paper,skipsamekey,11pt,english]{curve}

% 中文支持 (使用 XeLaTeX)
\usepackage[UTF8]{ctex}
%%%%%%%%%%%%%%%%%%%%%%%%%%%%%%%%%%%%%%
% 语言选择 / Language Selection
% ------------------------------------
% 启用下面一行以生成中文简历 / Uncomment for Chinese version:
\usepackage[Chinese]{languageSelection}
%
% 启用下面一行以生成英文简历 / Uncomment for English version:
% \usepackage[English]{languageSelection}
%%%%%%%%%%%%%%%%%%%%%%%%%%%%%%%%%%%%%%




% Default biblatex style used for the publication list is APA6. If you wish to use a different style or pass other options to biblatex you can change them here.

% 下面这行代码将全部作者显示缩写
% \PassOptionsToPackage{style=ieee,sorting=ydnt,uniquename=init,defernumbers=true}{biblatex}

% 下面这行代码将全部作者显示全称
\PassOptionsToPackage{style=ieee,sorting=none,uniquename=init,defernumbers=true,giveninits=false}{biblatex}

% Most commands and style definitions are in settings.sty.
\usepackage{settings}

% \DeclareSortingTemplate{firstme}{
%   \sort{
%     \field{firstauthor} % 优先根据这个字段：有 true 的会排前面
%   }
%   \sort{
%     \field{year}
%     \field{month}
%   }
%   \sort{
%     \field{author}
%   }
%   \sort{
%     \field{title}
%   }
% }
% If you need to further customise your biblatex setup e.g. with \DeclareFieldFormat etc please add them here AFTER loading settings.sty. For example, to remove the default "[Online] Available:" prefix before URLs when using the IEEE style:
\DefineBibliographyStrings{english}{url={\textsc{url}}}

%% Only needed if you want a Publication List
% \addbibresource{own-bib.bib}
\addbibresource{refs.bib}

%% Specify your last name(s) and first name(s) (as given in the .bib) to automatically bold your own name in the publications list. 
%% One caveat: You need to write \bibnamedelima where there's a space in your name for this to work properly; or write \bibnamedelimi if you use initials in the .bib
% \mynames{Lim/Lian\bibnamedelima Tze}

%% You can specify multiple names like this, especially if you have changed your name or if you need to highlight multiple authors. See items 6–9 in the example "Journal Articles" output.
% \mynames{Lim/Lian\bibnamedelima Tze,
%   Wong/Lian\bibnamedelima Tze,
%   Lim/Tracy,
%   Lim/L.\bibnamedelimi T.}
%% MAKE SURE THERE IS NO SPACE AFTER THE FINAL NAME IN YOUR \mynames LIST
% \mynames{Zhiqiang\bibnamedlima Ma}
\mynames{Ma/Zhiqiang}

% Change the fonts if you want
\ifxetexorluatex % If you're using XeLaTeX or LuaLaTeX
  \usepackage{fontspec} 
  %% You can use \setmainfont etc; I'm just using these font packages here because they provide OpenType fonts for use by XeLaTeX/LuaLaTeX anyway
  \usepackage[p,osf,swashQ]{cochineal}
  \usepackage[medium,bold]{cabin}
  \usepackage[varqu,varl,scale=0.9]{zi4}
\else % If you're using pdfLaTeX or latex
  \usepackage[T1]{fontenc}
  \usepackage[p,osf,swashQ]{cochineal}
  \usepackage{cabin}
  \usepackage[varqu,varl,scale=0.9]{zi4}
\fi

% Change the page margins if you want
% \geometry{left=1cm,right=1cm,top=1.5cm,bottom=1.5cm}

% Change the colours if you want
% \definecolor{SwishLineColour}{HTML}{00FFFF}
% \definecolor{MarkerColour}{HTML}{0000CC}

% Change the item prefix marker if you want
% \prefixmarker{$\diamond$}

% 设置 entry 日期（key）左对齐
\keyalignment{l}

%% Photo is only shown if "fullonly" is included
\includecomment{fullonly}
% \excludecomment{fullonly}


%%%%%%%%%%%%%%%%%%%%%%%%%%%%%%%%%%%%%%
% icos are from https://fontawesome.com/

\CN{\title{个人简历}}
\EN{\title{Curriculum Vitae}}

\CN{
\leftheader{%
{\LARGE\bfseries\sffamily 马志强，博士生} \\
\begin{tabular}{ll}
  
出生日期：19**年**月 &
籍$\quad \ \ \ $贯：** \\

政治面貌：**** &
电子邮箱：{\texttt{zhiqiang.ma97@gmail.com}}
   \\

\end{tabular}
}
}

\EN{
\leftheader{%
{\LARGE\bfseries\sffamily Zhiqiang Ma, Ph.D. Candidate} \\
\begin{tabular}{ll}
\makefield{\faEnvelope[regular]}{{\texttt{zhiqiang-ma@hotmail.com}}} &
  \makefield{\faEnvelope[regular]}{{\texttt{zhiqiang.ma97@gmail.com}}} \\
  % \makefield{\faEnvelope[regular]}{\href{mailto:example@gmail.com}{\texttt{zhiqiang-ma@hotmail.com}}}
  % fontawesome5 doesn't have the X icon so we use
  % the simpleicons package here instead; but some 
  % font size adjustment might be needed
  % \makefield{{\scriptsize\simpleicon{x}}}{\!\href{https://x.com/overleaf_example}{\texttt{@overleaf\_example}}}
  % \makefield{\faLinkedin}{\href{http://www.linkedin.com/in/example/}{\texttt{example}}}
\makefield{\faHome}{****, China} &
% \makefield{\faGithub}{}
\makefield{\faMars}{Male}
% \makefield{\faEarthAsia}{ss}

  %% Next line
  % \makefield{\faGlobe}{\url{http://example.example.org/}}
  % You can use a tabular here if you want to line up the fields.
\end{tabular}
}
}


\rightheader{~}
\begin{fullonly}
\photo[r]{photo}
% \photo[r]{photo}
\photoscale{0.13}
\end{fullonly}

% \title{Curriculum Vitae}

\begin{document}
\maketitle
\makeheaders[c]
\makerubric{short_biography}

\makerubric{education}
% \makerubric{employment}

% Example citation to ensure biblatex finds a citation
\nocite{*}

% If you're not a researcher nor an academic, you probably don't have any publications; delete this line.
%% Sometimes when a section can't be nicely modelled with the \entry[]... mechanism; hack our own and use \input NOT \makerubric
\CN{
    \makerubrichead{研究成果}
    \printbibliography[heading={subbibliography},title={会议论文},type=inproceedings]

    \printbibliography[heading={subbibliography},title={期刊论文},type=article]
}

\EN{
    %% Sometimes when a section can't be nicely modelled with the \entry[]... mechanism; hack our own
    \makerubrichead{Research Publications}
    % \textbf{\large Journal}
    % \subrubric{Journal}
    % \section*{Journal Articles}
    % \entry
    % {\cvlabelnum{1}} % 产生红圈编号
    % {\textbf{Z. Ma} and J. He, “Cryptanalysis and improvement of a multi-factor authenticated key exchange protocol,” \emph{International Journal of Network Security}, vol. 25, no. 5, pp. 764-776, Sep. 2023, ISSN: 1816-353X.}

    % \textbf{conference}
    % \subrubric{Conference}
    % \section*{Journal Articles}
    % \entry
    % {\cvlabelnum{1}}
    % {W. Jia, Z. Yang, and \textbf{Z. Ma}, “Lightweight leakage-resilient authenticated key exchange for industrial internet of things,” in \emph{TrustCom}, IEEE, 2024, pp. 1051-1059.}

    % \entry
    % {\cvlabelnum{2}}
    % {\textbf{Z. Ma} and J. He, “Outsider key compromise impersonation attack on a multi-factor authenticated key exchange protocol,” in \emph{ACNS Workshops}, ser. Lecture Notes in Computer Science, vol. 13285, Springer, 2022, pp. 320-337.}


    %% Assuming you've already given \addbibresource{own-bib.bib} in the main doc. Right? Right???
    % \nocite{*}

    %% If you just want everything in one list
    % \printbibliography[heading={none}]

    \printbibliography[heading={subbibliography},title={Conference},type=inproceedings]

    \printbibliography[heading={subbibliography},title={Journal},type=article]


    % \printbibliography[heading={subbibliography},title={Books and Chapters},filter={booksandchapters}]

}

\makerubric{skills}
\makerubric{misc}

% \makerubric{referee}
% \input{referee-full}

% \input{own-bib.bib}
% \input{refs.bib}

\end{document}